\begin{abstract}
本文针对脑机接口与精神性疾病诊断中的脑电图计算建模问题，围绕脑电信号的保真降噪、视觉诱发响应提取、神经生成机制建模及认知过程建模展开研究。

针对问题一，本文首先对Fz、F3、F4三通道脑电信号进行刺激事件对齐与基线校正，并构建“基础保真滤波--试次质量加权--候选响应子空间平滑”的分层预处理框架。为避免过度降噪破坏弱视觉响应，进一步采用已知波形注入、分块交叉验证、分半稳定性与预处理敏感性分析对候选方法进行约束选择；在此基础上提取不同实验项目、不同刺激条件下的事件相关电位，并采用低阶高斯基函数进行曲线参数化拟合。

针对问题二，\TODO{根据后续LGN--皮层--头皮生成模型补写：核心方程、关键参数、左右三角刺激差异的机制解释与验证结果。}

针对问题三，\TODO{根据后续认知过程模型补写：状态变量、模型结构、参数估计、实验验证与应用意义。}

最终，本文形成从“数据预处理--有效响应提取--神经机制解释--认知状态建模”的统一计算框架，为非侵入式脑机接口及神经精神疾病相关脑电研究提供可解释的建模路径。

\noindent\textbf{关键词：}脑电图；脑机接口；事件相关电位；保真降噪；神经计算模型；视觉认知
\end{abstract}
